\origpage{346}{N. Robertson, P.D. Seymour / Journal of Combinatorial Theory, Series B 90 (2004) 325--385}
\noindent $v \in \bar{x} \cap V(C)$. The atoms of $\Gamma^*$ corresponding to $\{v\}$ and to $x$ are adjacent in $\Gamma^*$, and so $d(x, \{v\}) \le 1$. But since $\{v\} \cap \Lambda = \emptyset$, it follows from the second assertion that $d(z, \{v\}) \ge \kappa$, and so
\begin{equation*}
d(z, x) \ge d(z, \{v\}) - d(x, \{v\}) \ge \kappa - 1.
\end{equation*}
This proves the third assertion.

To verify the fourth, let $f \in E(\Gamma)$ with $d(z, f) = rep(\mathscr{T})$. (This exists by Theorem 8.12 of [6].) Let $C^*$ be a circuit of $\Gamma^*$ with $U(C^*) \subseteq \bar{\Lambda}$ and with $|E(C^*)| < 2(rep(\mathscr{T}) - \kappa - 5)$. Let $D$ be the closed disc in $\bar{\Lambda}$ bounded by $U(C^*)$, and let $v \in V(C^*) \cap V(\Gamma)$. Then $d(z, v) \le \kappa + 5$, and so $d(v, f) \ge rep(\mathscr{T}) - \kappa - 5$. Hence $f \nsubseteq ins(C^*)$, because $|E(C^*)| < 2(rep(\mathscr{T}) - \kappa - 5)$, and so $ins(C^*)$ and $D$ are both closed discs in $\Sigma \setminus f$ bounded by $U(C^*)$. Consequently $D = ins(C^*)$, and so $ins(C^*) \subseteq \bar{\Lambda}$. This verifies the fourth assertion, and so completes the proof of 7.1. $\square$

\sectionhead{8. Tie-breakers}

Throughout this section $\Gamma$ is an internally 3-connected painting in $\Sigma$. Let $e \in E(\Gamma^*)$ with ends $v \in V(\Gamma)$ and $r^* \in R^*(\Gamma)$, and let $C$ be the circuit of $sk(\Gamma)$ with $U(C)$ bounding $r$. (We assume for the moment that $C$ exists.) Let $f_1, f_2$ be the two edges of $C$ incident with $v$, and let $c_1, c_2$ be the two edges of $\Gamma$ with $f_1 \subseteq c_1, f_2 \subseteq c_2$. Now if $(A, B)$ is a separation of $\Gamma$, we say that $(A, B)$ splits $e$ if $c_1 \in E(A)$ and $c_2 \in E(B)$ or vice versa. (If the circuit $C$ does not exist, then, since the perimeter of $r$ has no ``repeated'' vertices because $\Gamma$ is internally 3-connected, it follows that $|V(\Gamma)| = 2$ and $|E(\Gamma)| = 1$. In this case we say that $(A, B)$ does not split $e$.) We shall need the following lemma.

\textbf{8.1.} \textit{Let $(A, B)$ be a separation of $\Gamma$ of order $< rep(\Gamma)$, with $E(A), E(B) \ne \emptyset$. Then there is a circuit $F$ of $\Gamma^*$ such that $(A, B)$ splits every edge of $F$ and there is a disc $\Delta \subseteq \hat{\Sigma}$ bounded by $U(F)$ such that either $\Gamma \cap \Delta \subseteq A$ or $\Gamma \cap \Delta \subseteq B$.}

\textbf{Proof.} Let $G$ be the subdrawing of $\Gamma^*$ with $V(G) = V(\Gamma^*)$ and edges those edges of $\Gamma^*$ with are split by $(A, B)$. Now $E(G) \ne \emptyset$ since $E(A), E(B) \ne \emptyset$; and every vertex of $G$ has even degree from the definition of $G$. Moreover, for any circuit $F$ of $G$, half its vertices belong to $V(A \cap B)$, and so $\frac{1}{2}|E(F)| < rep(\Gamma)$; and therefore there is a disc in $\hat{\Sigma}$ bounded by $U(F)$. Let us choose a minimal disc $\Delta \subseteq \hat{\Sigma}$ with $bd(\Delta) \subseteq U(G)$; and let $F$ be the circuit of $\Gamma^*$ with $U(F) = bd(\Delta)$. Since every edge of $G$ is contained in a circuit of $G$, and $G \cap \Delta$ has no circuit except $F$, it follows that $G \cap \Delta = F$. But then $\Gamma \cap \Delta \subseteq A$ or $\Gamma \cap \Delta \subseteq B$, as required. $\square$

In [5] we discussed ``tie-breakers'' in general hypergraphs. In this paper we only need a particular kind of tie-breaker, chosen to work nicely in paintings, but we use the same name. Thus, for each $e \in E(\Gamma^*)$ let $\lambda(e) > 0$ be some real number, and for each $v \in V(\Gamma)$ let $\lambda(v) > 0$ be some real number, such that all the $\lambda(e)$'s and $\lambda(v)$'s are
